% =============================================================================
% PARTE III: Lo Specchio Digitale
% =============================================================================

\addcontentsline{toc}{part}{PARTE III: Lo Specchio Digitale} % Questo la aggiunge all'indice

\cleardoublepage
\thispagestyle{empty}

\vspace*{\fill}

\begin{center}
    % Numero della parte
    {\sffamily\huge\textls[100]{\MakeUppercase{Parte III}}}
    
    \vspace{1.5em}
    
    % Linea decorativa
    \rule{0.5\textwidth}{0.4pt}
    
    \vspace{1.5em}
    
    % Titolo della parte
    {\sffamily\Huge\bfseries Lo Specchio Digitale}
    
    \vspace{2em}
    
    % Sottotitolo/descrizione
    \begin{minipage}{0.7\textwidth}
        \centering
        {\large\itshape\color{moderngray}
        Come l'Intelligenza Artificiale riflette,\\
        amplifica e cambia la nostra mente.}
    \end{minipage}
    
    \vspace{1.5em}
    
    % Linea decorativa
    \rule{0.5\textwidth}{0.4pt}
\end{center}

\vspace*{\fill}

\clearpage
\thispagestyle{empty}
\vspace*{2em}
\begin{center}
	{\sffamily\large\bfseries\color{darkgray} Dal mondo alle macchine}
\end{center}
\vspace{1em}

La Parte II ci ha mostrato una mente che decide fuori posto. Le scorciatoie che funzionano nella vita di ogni giorno diventano errori sistematici nei mercati, nelle corsie, nei tribunali, davanti a un rischio lontano. Le cornici culturali fissano in anticipo che cosa ci sembra normale. L'infanzia e la famiglia le trasmettono, e il successo dei pochi che ce l'hanno fatta le conferma. Infine abbiamo scoperto che i nostri archivi somigliano agli aerei che tornavano alla base, conservano soltanto ciò che qualcuno ha visto, registrato e scritto in un certo modo.

La domanda di questa parte nasce da qui. Che cosa impara una macchina da un archivio così? E che cosa succede quando quella macchina comincia a restituirci il mondo, scegliendo che cosa vediamo, rispondendoci con voce umana, affiancando il nostro giudizio e producendo a sua volta immagini, testi e voci?

Qui il titolo del libro diventa letterale. Un sistema di intelligenza artificiale è una memoria del passato messa al lavoro per prevedere il futuro. Se la memoria è distorta, la previsione eredita la distorsione e la proietta in avanti, con l'autorevolezza dei numeri. Il problema non è soltanto nella macchina, e non è soltanto in noi. Vedremo che il bias può nascere nei dati, nell'obiettivo, nella progettazione, nell'interazione e in chi legge la risposta.

Senza questa parte la Parte IV sarebbe una raccolta di buoni consigli. Per orientarsi bisogna prima sapere dove, esattamente, il nostro modo di sbagliare e quello delle macchine si incontrano.

\cleardoublepage
